% SPDX-FileCopyrightInfo: Copyright © DuMux Project contributors, see AUTHORS.md in root folder
% SPDX-License-Identifier: CC-BY-4.0

\documentclass[11pt,a4paper,headinclude,footinclude,DIV16]{scrreprt}
\usepackage[usenames,dvipsnames,table]{xcolor} % include at beginning because of compitibility issues
\usepackage[english]{babel}
\usepackage[utf8]{inputenc}
\usepackage[automark]{scrlayer-scrpage}
\usepackage[normalem]{ulem}
\usepackage{amsmath}
\usepackage{amsfonts}
\usepackage{amssymb}
\usepackage{booktabs}
\usepackage{enumerate}
\usepackage{hyperref}
\usepackage{graphicx}
\usepackage{mhchem}

\usepackage{listings}
\usepackage[square,numbers]{natbib}
\usepackage{pdflscape}
\usepackage{rotating}
\usepackage{subfig}
\usepackage{tabularx}
\usepackage{theorem}
\usepackage{tikz}
\usepackage{tikz-qtree}
\usepackage{units}
\usepackage{url}
\usepackage{xspace}
\usepackage{accsupp}
\hypersetup{bookmarksdepth=3}
\usetikzlibrary{arrows}
\usetikzlibrary{backgrounds}
\usetikzlibrary{decorations.pathmorphing}
\usetikzlibrary{fit}
\usetikzlibrary{patterns}
\usetikzlibrary{positioning}
\usetikzlibrary{shapes}
\usetikzlibrary{trees}
\newcommand{\textcolordummy}[2]{#2}
\newcommand{\textcolordigits}{\textcolor{Mulberry}}
\lstset{showstringspaces=false,
 inputencoding=ansinew,
 breaklines=true,
 tabsize=4,
 columns=flexible,
 literate={0}{{\textcolordigits{0}}}{1}%
          {1}{{\textcolordigits{1}}}{1}%
          {2}{{\textcolordigits{2}}}{1}%
          {3}{{\textcolordigits{3}}}{1}%
          {4}{{\textcolordigits{4}}}{1}%
          {5}{{\textcolordigits{5}}}{1}%
          {6}{{\textcolordigits{6}}}{1}%
          {7}{{\textcolordigits{7}}}{1}%
          {8}{{\textcolordigits{8}}}{1}%
          {9}{{\textcolordigits{9}}}{1}%
          {.0}{{\textcolordigits{.0}}}{2}% Following is to ensure that only periods
          {.1}{{\textcolordigits{.1}}}{2}% followed by a digit are changed.
          {.2}{{\textcolordigits{.2}}}{2}%
          {.3}{{\textcolordigits{.3}}}{2}%
          {.4}{{\textcolordigits{.4}}}{2}%
          {.5}{{\textcolordigits{.5}}}{2}%
          {.6}{{\textcolordigits{.6}}}{2}%
          {.7}{{\textcolordigits{.7}}}{2}%
          {.8}{{\textcolordigits{.8}}}{2}%
          {.9}{{\textcolordigits{.9}}}{2}%
          {e0}{{\textcolordigits{e0}}}{2}% Following is to ensure that only periods
          {e1}{{\textcolordigits{e1}}}{2}% followed by a digit are changed.
          {e2}{{\textcolordigits{e2}}}{2}%
          {e3}{{\textcolordigits{e3}}}{2}%
          {e4}{{\textcolordigits{e4}}}{2}%
          {e5}{{\textcolordigits{e5}}}{2}%
          {e6}{{\textcolordigits{e6}}}{2}%
          {e7}{{\textcolordigits{e7}}}{2}%
          {e8}{{\textcolordigits{e8}}}{2}%
          {e9}{{\textcolordigits{e9}}}{2}%
          {e-0}{{\textcolordigits{e-0}}}{2}% Following is to ensure that only periods
          {e-1}{{\textcolordigits{e-1}}}{2}% followed by a digit are changed.
          {e-2}{{\textcolordigits{e-2}}}{2}%
          {e-3}{{\textcolordigits{e-3}}}{2}%
          {e-4}{{\textcolordigits{e-4}}}{2}%
          {e-5}{{\textcolordigits{e-5}}}{2}%
          {e-6}{{\textcolordigits{e-6}}}{2}%
          {e-7}{{\textcolordigits{e-7}}}{2}%
          {e-8}{{\textcolordigits{e-8}}}{2}%
          {e-9}{{\textcolordigits{e-9}}}{2},
}

% for listings of bash code
\lstdefinestyle{Bash}{
 language=Bash,
 backgroundcolor=\color{lightgray},
 basicstyle=\ttfamily\small\let\textcolor\textcolordummy,
 numbers=none,
 captionpos=b,
 tabsize=4,
 frame=single,
 rulecolor=\color{lightgray},
 framerule=1pt,
 framesep=1pt,
 rulesep=0pt,
 aboveskip=\bigskipamount,
 belowskip=\bigskipamount
}

% for listings of shell code
\lstdefinestyle{Shell}{
 language=sh,
 numbers=left,
 numbersep=5pt,
 numberstyle=\tiny\noncopynumber,
 basicstyle=\ttfamily\scriptsize,
 stringstyle=\color{BrickRed}\ttfamily\let\textcolor\textcolordummy,
 commentstyle=\color[gray]{0.35}\ttfamily\itshape\let\textcolor\textcolordummy,
}

% for listings of DuMuX code
\lstdefinestyle{DumuxCode}{
  language=C++,
  extendedchars=true,
  escapeinside={/*@}{@*/}, % escape the long comments
  numbers=left,
  numbersep=5pt,
  numberstyle=\tiny\noncopynumber,
  basicstyle=\ttfamily\scriptsize,
  keywordstyle=\color{dumuxBlue}\ttfamily\let\textcolor\textcolordummy,
  stringstyle=\color{BrickRed}\ttfamily\let\textcolor\textcolordummy,
  commentstyle=\color[gray]{0.35}\ttfamily\itshape\let\textcolor\textcolordummy,
  emph={NEW_TYPE_TAG, NEW_PROP_TAG, UNSET_PROP, TTAG, PTAG,
       SET_PROP, GET_PROP, GET_PROP_VALUE, GET_PROP_TYPE,
       SET_BOOL_PROP, SET_STRING_PROP, SET_SCALAR_PROP, SET_TYPE_PROP, SET_INT_PROP,
       GET_BOOL_PROP, GET_STRING_PROP, GET_SCALAR_PROP, GET_TYPE_PROP, GET_INT_PROP,
       GET_PARAM, GET_PARAM_FROM_GROUP, GET_RUNTIME_PARAM, GET_RUNTIME_PARAM_FROM_GROUP},
  emphstyle=\color{OliveGreen}\let\textcolor\textcolordummy,
  directivestyle=\color{Brown}\ttfamily\let\textcolor\textcolordummy,
}

% for listings of DuMuX parameter files
\lstdefinestyle{DumuxParameterFile}{
 language=sh,
 numbers=left,
 numbersep=5pt,
 numberstyle=\tiny\noncopynumber,
 basicstyle=\ttfamily\scriptsize,
 stringstyle=\color{BrickRed}\ttfamily\let\textcolor\textcolordummy,
 commentstyle=\color[gray]{0.35}\ttfamily\itshape\let\textcolor\textcolordummy,
 morecomment=[l][\color{dumuxBlue}\let\textcolor\textcolordummy]{[},
}

\newcommand{\noncopynumber}[1]{
    \BeginAccSupp{method=escape,ActualText={}}
    #1
    \EndAccSupp{}
}


\DeclareGraphicsExtensions{.pdf, .jpg}

% Dune and Dumux logo
\newcommand{\Dune}{{DUNE}\xspace}
\newcommand{\Dumux}{\texorpdfstring{Du\-Mu$^\text{x}$\xspace}{DuMuX\xspace}}
\newcommand{\DumuxVersion}{3.7}
\newcommand{\DumuxOldVersion}{3.6}
\definecolor{dumuxYellow}{HTML}{E19417}
\definecolor{dumuxBlue}{HTML}{0C73CF}
\newcommand{\DuneVersion}{2.9}
\newcommand{\DuneOldVersion}{2.8}
\newcommand{\CMakeVersion}{3.14}

% styles
\newcommand{\nextline}{\par\phantom{a}\vspace*{0.1\textwidth}}
\newcommand{\snakeline}{\uwave{\mbox{}}}
\DeclareRobustCommand\Cplusplus{\texorpdfstring{C\nolinebreak[4]\hspace{-.05em}\raisebox{.4ex}{\tiny\bfseries ++}\xspace}{C++}}

% notation
\newcommand{\porosity}{\phi}
\newcommand{\saturation}{S}

% a new counter you can give a label to it and thus reference it
% syntax: \numberThis{printedTextToBeLabeled}{label}
% if you wanted a \newline after a numbered thing, you could just add a empty line after ``\label{#2}''
\newcounter{thingCounter}
\renewcommand{\thethingCounter}{\arabic{thingCounter}}
\newcommand{\numberThis}[2]{%
        \refstepcounter{thingCounter}%
        \thethingCounter.\ #1 \label{#2}
}

%The theorems
\theorembodyfont{\upshape}
\theoremheaderfont{\sffamily\bfseries}
\newtheorem{lst}{Listing}

\DeclareMathOperator{\grad}{\mathbf{grad}}
\DeclareMathOperator{\curl}{curl}
\DeclareMathOperator{\Div}{div}
\newcommand{\meas}[1]{\lvert{#1}\rvert}

\pagestyle{scrheadings}

\title{
\begin{center}
\includegraphics[width=0.7\textwidth]{../logo/dumux_logo_hires_whitebg.png}
\\[3cm]
{\Huge Handbook}
\end{center}
}

\author{}

\date{Version \DumuxVersion, \\
Handbook version from \today}

\publishers{%
\vspace{10mm}
{\normalsize Lehrstuhl f\"ur Hydromechanik und Hydrosystemmodellierung, \\
Universit\"at Stuttgart, Paffenwaldring 61, D-70569 Stuttgart, Germany}\\
%
\bigskip
{\normalsize \texttt{\url{http://dumux.org}}}\\
}

\begin{document}

\maketitle

\setcounter{tocdepth}{1}
\tableofcontents
\newpage

\chapter{Overview and Infrastructure}
The handbook is in the process of being dissolved.

% SPDX-FileCopyrightInfo: Copyright © DuMux Project contributors, see AUTHORS.md in root folder
% SPDX-License-Identifier: CC-BY-4.0

\section{Directory Structure}
\label{sc_structure}

\Dumux has the following folder structure, which is similar to other \Dune modules.
\begin{itemize}
\item \texttt{bin}: binaries, e.g. used for the automatic testing, post-processing, installation
\item \texttt{cmake}: the configuration options for building \Dumux
\item \texttt{doc}: files necessary for the Doxygen documentation and this handbook, and various logos
\item \texttt{dumux}: the main folder, containing the source files. See Fig. \ref{fig:dumux-structure}
      for a visualized structure. For more information on the models, have a look at the
      Doxygen documentation.
\item \texttt{examples}: well-documented examples of applying \Dumux to typical physical problems. In the \texttt{README.md} files, the setup is explained, the used code is presented as well as documented and images resulting from the simulation are included. The \texttt{README.md} files are located directly in the subfolders of \texttt{examples} and can be displayed by web browsers.
\item \texttt{test}: tests for each numerical model and some functionality.
      The structure is equivalent to the \texttt{dumux} folder, the \texttt{references} folder
      contains solutions for the automatic testing. Each test program consist of a main file
      \texttt{main.cc}, the problem definition \texttt{*problem.hh} (specifying initial and boundary
      conditions), and an input file \texttt{params.input}.
      If necessary, spatially dependent parameters are defined in \texttt{*spatialparameters.hh}.
      For more detailed descriptions of the tests, please have a look at the Doxygen documentation.
\end{itemize}

\begin{figure}
% \begin{sidewaysfigure}
\begin{tikzpicture}[scale=0.8,grow'=right,level distance=1.5in,sibling distance=.05in]
\tikzset{edge from parent/.style={thick, draw=dumuxBlue, edge from parent fork right}}
\tikzset{every tree node/.style={draw, thick, align=center}}
\tikzset{frontier/.style={distance from root=6.0in}}

\tikzset{FirstLevel/.style={
  draw=dumuxYellow,
  rectangle,
  align=center,
  minimum width=1.1in,
  minimum height=0.2in,
  text width=1.2in,
}}
\tikzset{SecondLevel/.style={
  draw=dumuxBlue,
  rectangle,
  align=center,
  minimum width=1.1in,
  minimum height=0.2in,
  text width=1.2in,
}}

\tikzset{ThirdLevel/.style={
  draw=none,
  align=left,
  minimum width=4.2in,
  text width=4.1in,
  font=\footnotesize
}}


\Tree
[.\node[draw=dumuxYellow, ultra thick] {dumux};
  [.\node[FirstLevel] {adaptive};
    \node[ThirdLevel] {
          Data transfer on new grid, adaptation indicators.};
  ]
  [.\node[FirstLevel] {assembly};
    \node[ThirdLevel] {
          Matrix assembler and residual for all discretization methods.};
  ]
  [.\node[FirstLevel] {common};
    \node[ThirdLevel] {
          Common files of all models:
          definition of boundary conditions, time stepping, splines, dimensionless numbers ...};
    [.\node[SecondLevel] {geometry};
      \node[ThirdLevel] {Geometrical algorithms};
    ]
    [.\node[SecondLevel] {properties};
      \node[ThirdLevel] {Base properties for all models.};
    ]
    [.\node[SecondLevel] {typetraits};
      \node[ThirdLevel] {Helper classes to query type information on compile-time. };
    ]
  ]
  [.\node[FirstLevel] {discretization};
      \node[ThirdLevel] {Common methods for all discretizations (box, cell-centered TPFA/MPFA, staggered grid): variable caching, advective and diffusive fluxes, ...};
  ]
  [.\node[FirstLevel] {flux};
    [\node[ThirdLevel] {
            Collection of classes used to calculate advective and diffusive fluxes.};
    ]
  ]
  [.\node[FirstLevel] {freeflow};
    [.\node[SecondLevel] {\emph{models}};
      \node[ThirdLevel] {Single-phase free flow models using Navier-Stokes
                         and eddy-viscosity based Reynolds-averaged Navier-Stokes turbulence models, and shallow water equation model.};
    ]
  ]
  [.\node[FirstLevel] {geomechanics};
    [.\node[SecondLevel] {\emph{models}};
      \node[ThirdLevel] {Elastic and poro-elastic geomechanics models.};
    ]
  ]
  [.\node[FirstLevel] {io};
    \node[ThirdLevel] {Additional in-/output possibilities like restart files, gnuplot-interface,
                       VTKWriter extensions and files for grid generation.};
  ]
  [.\node[FirstLevel] {linear};
    \node[ThirdLevel] {Linear solver backend.};
  ]
  [.\node[FirstLevel] {material};
    [.\node[SecondLevel] {binarycoefficients};
      \node[ThirdLevel] {Binary coefficients (like binary diffusion coefficients) and those
                         needed for the constitutive relationships (e.g. Henry coefficient)};
    ]
    [.\node[SecondLevel] {chemistry};
      \node[ThirdLevel] {Files needed to account for, e.g. electro-chemical processes as in a fuel cell.};
    ]
    [.\node[SecondLevel] {components};
      \node[ThirdLevel] {Properties of a pure chemical substance (e.g. water)
                         or pseudo substance (e.g. air).};
    ]
    [.\node[SecondLevel] {constraintsolvers};
      \node[ThirdLevel] {Constraint solvers to make sure that the resulting fluid state is consistent with a
                         given set of thermodynamic equations.};
    ]
    [.\node[SecondLevel] {eos};
      \node[ThirdLevel] {Equations of state (eos) are auxiliary classes which provide
                         relations between a fluid phase's temperature, pressure, composition
                         and density.};
    ]
    [.\node[SecondLevel] {fluidmatrixint.};
      \node[ThirdLevel] {Constitutive relationships (e.g. capillary pressures, relative permeabilities)};
    ]
    [.\node[SecondLevel] {fluidstates};
      \node[ThirdLevel] {Fluid states are responsible for caching the thermodynamic
                         configuration of a fluid system at a given spatial and temporal position.};
    ]
    [.\node[SecondLevel] {fluidsystems};
      \node[ThirdLevel] {Fluid systems express the thermodynamic relations between quantities.};
    ]
    [.\node[SecondLevel] {solidstates};
      \node[ThirdLevel] {Solid states are responsible for caching the thermodynamic
                         configuration of a solid system at a given spatial and temporal position.};
    ]
    [.\node[SecondLevel] {solidsystems};
      \node[ThirdLevel] {Solid systems express the thermodynamic properties of a solid.};
    ]
    [.\node[SecondLevel] {spatialparams};
      \node[ThirdLevel] {Base class for all spatially dependent variables, like permeability and
                         porosity. Includes spatial averaging routines. All other properties are
                         specified in the specific files of the respective models.};
    ]
  ]
  [.\node[FirstLevel] {multidomain};
    \node[ThirdLevel] {
          Common infrastructure to couple multiple domains, models or physics.};
    [.\node[SecondLevel] {embedded};
      \node[ThirdLevel] {Embedding a lower-dimensional model into a higher-dimensional one};
    ]
    [.\node[SecondLevel] {facet};
      \node[ThirdLevel] {Mixed-dimensional coupling at facets.};
    ]
    [.\node[SecondLevel] {boundary};
      \node[ThirdLevel] {Coupling at the domain boundaries.};
    ]
  ]
  [.\node[FirstLevel] {nonlinear};
    \node[ThirdLevel] {Newton's method.};
  ]
  [.\node[FirstLevel] {parallel};
    \node[ThirdLevel] {Helper files for parallel simulations.};
  ]
  [.\node[FirstLevel] {porousmediumflow};
    [.\node[SecondLevel] {\emph{models}};
    \node[ThirdLevel] {Specific definition for porous medium flow simulations for all models:
                       implementation of equations,
                       model specific properties and indices.};
    ]
  ]
]
\end{tikzpicture}
\caption{Structure of the directory \texttt{dumux} containing the \Dumux source files.}
\label{fig:dumux-structure}
% \end{sidewaysfigure}
\end{figure}

% SPDX-FileCopyrightInfo: Copyright © DuMux Project contributors, see AUTHORS.md in root folder
% SPDX-License-Identifier: CC-BY-4.0

\section{Setup of new Folders and new Tests}
\label{sc_newfoldersetup}
This section describes how to set up a new folder and how to tell
the build system there is a new one.
\paragraph{Adding new Folders}
\begin{enumerate}[1)]
 \item create new folder with content
 \item adapt the \verb+CMakeList.txt+ in the folder above by adding a line with
       \verb+add_subdirectory(NEW_FOLDER)+
 \item create a \verb+CMakeList.txt+ in the newly created folder
 \item go to your \texttt{build}-directory and type \verb+make+ to
       re-configure the system
\end{enumerate}

\paragraph{Adding new Test Programs}
\noindent To add a test use the \texttt{add\_dune\_test} macro within the \texttt{CMakeList.txt} file.
The macro can be used with a variable amount of arguments. A simple call could look like this:

\begin{lstlisting}[style=DumuxCode]
dumux_add_test(NAME my_test
               SOURCES main.cc
               CMD_ARGS my_test params.input)
\end{lstlisting}

Here, we create an executable called \texttt{my\_test} from a source file \texttt{main.cc}.
The name of the test will also be \texttt{my\_test} (has to be unique). The last argument specifies a command - here, we just run the executable \texttt{my\_test} with an input file \texttt{params.input}. For more advanced uses of
the \texttt{add\_dune\_test} macro, have a look at the \texttt{test} directory. A complete documentation is given under \url{https://www.dune-project.org/sphinx/core-2.7/}.

\section{Scripts}
\label{sc_scripts}

A suite of scripts is available within the dumux repository to assist in the performance of common tasks. These scripts can be found in the \texttt{dumux/bin/} directory. 
The majority of these scripts began as bash (*.sh) scripts, but most have been replaced with more portable and user-friendly Python versions. 
Some of these scripts are used in the installation or the creation of working dumux-dependent repositorities; these are outlined in section \ref{ssec:repoDevScripts}.
Others can be used to perform post-processing tasks after simulation data has been produced; these are outlined in section \ref{ssec:postprocessingScripts}. In order to maintain the suite of unit tests and examples within this repository, some testing scripts are also required; these are outlined in section \ref{ssec:testingScripts}. All remaining tests are mentioned in section \ref{ssec:miscScripts}.

\subsection{Repository development Scripts:} \label{ssec:repoDevScripts}

\begin{itemize}
  \item Installation Scripts: \texttt{installdumux.py} and \texttt{installexternal.py} both have been developed to make the installation of the base Dumux suite and any external dependencies as simple as possible. 
  \item Module Development: When developing within the dumux simulation environment, it is recommended to do so within your own module. When sharing this module with others, or installing it on multiple machines, it is often difficult to find exactly which version you are using so that your work will run in different places. To simplify this, the following scripts are available in the \texttt{util/} directory: 
  	\begin{itemize}
  		\item The script \texttt{createdockerimage.py} can help to create a docker image to put your work in a container for simple delivery.
		\item In order to extract the content and dependencies within your model, \texttt{extractmodulepart.py} will collect all headers used to run a specified test and collect them all in a functioning external module.
		\item To document the exact versions that have been used, the scripts \texttt{getusedversions.py} is available to search each through an existing simulation dumux suite and find the branch, commit hash, and commit date of the current implemention. 
		\item The script \texttt{create\_cmakelists.py} can be used to integrate new header files to the dune/dumux cmake system.
		\item In order to produce an installscript that will install your repository and its dependencies on another computer, the \texttt{makeinstallscript.py} is available. This script will collect all of the dependencies and their versions using the \texttt{getmoduleinfo.py} and the other scripts listed above and write a script to install a new version in the same configuration. In the past it has been recommended to run the script \texttt{test\_dumux.sh} to ensure that the final dunecontrol and compilation pass have worked. 
  	\end{itemize} 
\end{itemize}

\subsection{Post-processing Scripts:} \label{ssec:postprocessingScripts}
Data visualization and post processing is an integral part of any exercise in simulation. A few scripts are available to begin with in the \texttt{postprocessing/} subdirectory. 
\begin{itemize}
	\item Paraview data collection scripts: Although there are many tools available within the open-source visualization program paraview (see \ref{ssec:paraview}), performing consistent data collection from your simulations can be tedious when analyzing many output files. The scripts \texttt{exportscreenshot2d.py}, \texttt{extractlinedata.py}, and \texttt{extractpointdataovertime.py} are each scripts that can help to perform paraview data collection commands consistently and each produce easy to use \texttt{.csv} files or images. 
	\item Error convergence script: When evaluting errors to a reference solution across grid refinements, the L2 error is typically evaluated. A script to help perform this is available in \texttt{l2error.py}.
\end{itemize}

\subsection{Testing Scripts:} \label{ssec:testingScripts}
In order to ensure consistency when developing in dumux, developers are encouraged to introduce unit tests and reference tests that can be run to ensure new code introductions do not break existing functions. While existing simulations should each test something when the \texttt{ctest} command is called, a few scripts exist in the \texttt{testing/} subdirectory to make this simple to develop and quick to evaluate.
\begin{itemize}
	\item Run tests: In order to automatically find and run tests that fit a specific critera, the scripts \texttt{findtests.py}, \texttt{runtest.py} and \texttt{runselectedtests.py} each are available. These are each used to automate the testing process.  
	\item Reference Solution Comparisons: The scripts \texttt{fuzzycomparevtu.py} and \texttt{fuzzycomparedata.py} are both used to compare simulation results against reference simulation results. When running a test, one can use these files to compare the output from the test against reference simulation results. Should these results differ, the difference for each field is shown. These scripts are typically directly called during testing. 
\end{itemize}

\subsection{Miscellaneous:} \label{ssec:miscScripts}
In addition, a few miscellaneous scripts are available.
\begin{itemize}
	\item Parameter scripts: The script \texttt{doc/generate_parameterlist.py} is available for tracking all of the runtime parameters used.
	\item Other: \texttt{doc/getcontributors.sh} can be used to find the names and contact information for all commit authors over a given time period. 
\end{itemize}

\section{Assembling the linear system}
\label{sc_linearsystem}
The physical system is implemented as the mathematical differential equation in
local operators. \Dumux generates the linear system automatically. Read on, to
learn what is done internally.

\subsection{Newton's method}
The differential equations are implemented in the residual form. All terms are
on the left hand side and are summed up. The terms contain values for the primary
variables which are part of the solution vector $\textbf{u}$. The sum of the terms
is called residual $\textbf{r}(\textbf{u})$ which is a function of the solution. For
example:
\begin{align*}
\underbrace{
  \phi \frac{\partial \varrho_\alpha S_\alpha}{\partial t}
 -
 \text{div} \left(
 \varrho_\alpha \frac{k_{r\alpha}}{\mu_\alpha} \mathbf{K}
 \left(\grad\, p_\alpha - \varrho_{\alpha} \mathbf{g} \right)
 \right) - q_\alpha} _
{=: \, \textbf{r}(\textbf{u})}
= 0
\end{align*}

We don't know the solution $\textbf{u}$, so we use the iterative Newton's method to
obtain a good estimate of $\textbf{u}$. We start with an initial guess $\textbf{u}^0$ and
calculate it's residual $\textbf{r}(\textbf{u}^0)$. To minimize the error, we calculate
the derivative of the residual with respect to the solution. This is the Jacobian
matrix
\begin{align*}
  \frac{\text{d}}{\text{d}\textbf{u}}\textbf{r} \left(\textbf{u}^i\right)
  = J_{\textbf{r} \left(\textbf{u}^i\right)}
  = \left(\frac{\text{d}}{\text{d}\textbf{u}^i_m}\textbf{r} \left(\textbf{u}^i\right)_n\right)_{m,n}
\end{align*}
with $i$ denoting the Newton iteration step.
Each column is the residual derived with respect to the $m$th entry of $\textbf{u}^i$.

The Jacobian indicates the direction where the residual increases. By solving the
linear system
\begin{align*}
  J_{\textbf{r}(\textbf{u}^i)} \cdot \textbf{x}^i = \textbf{r}(\textbf{u}^i)
\end{align*}
we calculate the direction of maximum growth $\textbf{x}^i$. We subtract it from
our current solution to get a new, better solution
$\textbf{u}^{i+1} = \textbf{u}^i - \textbf{x}^i$.

We repeat the calculation of of the Jacobian $J_{\textbf{r}(\textbf{u}^i)}$ and the
direction of maximum growth $\textbf{x}^i$ until our approximated solution becomes good enough.

\subsection{Structure of matrix and vectors}
To understand the meaning of an entry in the matrix or the vector of the linear system, we have
to define their structure. Both have a blocking structure. Each block contains the degrees of
freedom (also called variable or unknown) for a control volume. The equation index is used
to order of the degrees of freedom. For each control volume we have one block. The mapper is
used to order the blocks.

\begin{figure}[htbp]
\begin{center}
\begin{tikzpicture}[fill=dumuxBlue]
  %% blocking structure
  % matrix
  \node at (0.3,4.2){\footnotesize 1. CV};
  \node at (1.7,4.2){\footnotesize 2. CV};
  \node at (3.5,4.2){\footnotesize $n$. CV};

  \draw (0,0) rectangle (4,4);

  \fill (0.1,3.1) rectangle (0.9,3.9);
  \fill (1.1,3.1) rectangle (1.9,3.9);
  \node at (2.5,3.5) {$\dots$};
  \fill (3.1,3.1) rectangle (3.9,3.9);
  \node at (4,3.5) [right]{\footnotesize 1. CV};

  \fill (0.1,2.1) rectangle (0.9,2.9);
  \fill (1.1,2.1) rectangle (1.9,2.9);
  \node at (2.5,2.5) {$\dots$};
  \fill (3.1,2.1) rectangle (3.9,2.9);
  \node at (4,2.5) [right]{\footnotesize 2. CV};

  \node at (0.5,1.5) {$\vdots$};
  \node at (1.5,1.5) {$\vdots$};
  \node at (2.5,1.5) {$\ddots$};
  \node at (3.5,1.5) {$\vdots$};

  \fill (0.1,0.1) rectangle (0.9,0.9);
  \fill (1.1,0.1) rectangle (1.9,0.9);
  \node at (2.5,0.5) {$\dots$};
  \fill (3.1,0.1) rectangle (3.9,0.9);
  \node at (4,0.5) [right]{\footnotesize $n$. CV};

  % vector
  \draw (5.5,0) rectangle (5.9,4);
  \fill (5.6,3.1) rectangle (5.8,3.9);
  \fill (5.6,2.1) rectangle (5.8,2.9);
  \node at (5.7,1.5) {$\vdots$};
  \fill (5.6,0.1) rectangle (5.8,0.9);

  %% intra-block structure
  \fill (8.1,2.1) rectangle (8.9,2.9);
  \draw (9,2.8) -- (9.6,3.4);
  \draw (9,2.6) -- (9.6,2.8);
  \draw (9,2.2) -- (9.3,1.6);

  \node at (10,4) {${eqIdx}$};
  \node at (10,3.4) {$0$};
  \node at (10,2.8) {$1$};
  \node at (10,2.2) {$\vdots$};
  \node at (10,1.6) {$m-1$};

  \fill (11.1,2.1) rectangle (11.3,2.9);
  \draw (11,2.8) -- (10.4,3.4);
  \draw (11,2.6) -- (10.4,2.8);
  \draw (11,2.2) -- (10.7,1.6);
\end{tikzpicture}
\end{center}
\caption{Structure of matrix and vector, left blocking structure, right within block}
\end{figure}

Accessing entries follows this structure. You can access the pressure value in the third sub-control volume in
a vector \lstinline{sol} with \lstinline{sol[2][pressureIdx]}.


\chapter{Advanced \Dumux\ -- Detailed Instructions}
This chapter contains detailed information for those who are interested
in deeper modifications of underlying \Dumux models, classes, functions, etc.
% SPDX-FileCopyrightInfo: Copyright © DuMux Project contributors, see AUTHORS.md in root folder
% SPDX-License-Identifier: CC-BY-4.0

\section{Temporal Discretization and Solution Strategies}
In this section, the temporal discretization as well as solution strategies (monolithic/sequential) are presented.

\subsection{Temporal discretization}

Our systems of partial differential equations are discretized in space and in time.

Let us consider the general case of a balance equation of the following form
\begin{equation}\label{eq:generalbalance}
\frac{\partial m(u)}{\partial t} + \nabla\cdot\mathbf{f}(u, \nabla u) + q(u) = 0,
\end{equation}
seeking an unknown quantity $u$ in terms of storage $m$, flux $\mathbf{f}$ and source $q$.
All available Dumux models can be written mathematically in form of \eqref{eq:generalbalance}
with possibly vector-valued quantities $u$, $m$, $q$ and a tensor-valued flux $\mathbf{f}$.
For the sake of simplicity, we assume scalar quantities $u$, $m$, $q$ and a vector-valued
flux $\mathbf{f}$ in the notation below.

For discretizing \eqref{eq:generalbalance}, we need to choose an
approximation for the temporal derivative $\partial m(u)/\partial t$.
While many elaborate methods for this approximation exist,
we focus on the simplest one of a first order difference quotient
\begin{equation}\label{eq:euler}
\frac{\partial m(u_{k/k+1})}{\partial t}
\approx \frac{m(u_{k+1}) - m(u_k)}{\Delta t_{k+1}}
\end{equation}
for approximating the solution $u$ at time $t_k$ (forward) or $t_{k+1}$ (backward).
The question of whether to choose the forward or the backward quotient leads to the
explicit and implicit Euler method, respectively.
In case of the former, inserting \eqref{eq:euler} in \eqref{eq:generalbalance}
at time $t_k$ leads to
\begin{equation}\label{eq:expliciteuler}
\frac{m(u_{k+1}) - m(u_k)}{\Delta t_{k+1}} + \nabla\cdot\mathbf{f}(u_k, \nabla u_k) + q(u_k) = 0,
\end{equation}
whereas the implicit Euler method is described as
\begin{equation}\label{eq:impliciteuler}
\frac{m(u_{k+1}) - m(u_k)}{\Delta t_{k+1}}
+ \nabla\cdot\mathbf{f}(u_{k+1}, \nabla u_{k+1}) + q(u_{k+1}) = 0.
\end{equation}
Once the solution $u_k$ at time $t_k$ is known, it is straightforward
to determine $m(u_{k+1})$ from \eqref{eq:expliciteuler},
while attempting to do the same based on \eqref{eq:impliciteuler}
involves the solution of a system of equations.
On the other hand, the explicit method \eqref{eq:expliciteuler} is stable only
if the time step size $\Delta t_{k+1}$ is below a certain limit that depends
on the specific balance equation, whereas the implicit method \eqref{eq:impliciteuler}
is unconditionally stable.

\subsection{Solution strategies to solve equations}
The governing equations of each model can be solved monolithically or sequentially.
The basic idea of the sequential algorithm is to reformulate the
equations of multi-phase flow into one equation for
pressure and equations for phase/component/... transport. The pressure equation
is the sum of the mass balance equations and thus considers the total flow of the
fluid system. The new set of equations is considered as decoupled (or weakly coupled)
and can thus be solved sequentially. The most popular sequential model is the
fractional flow formulation for two-phase flow which is usually implemented applying
an IMplicit Pressure Explicit Saturation algorithm (IMPES).
In comparison to solving the equations monolithically, the sequential structure allows the use of
different discretization methods for the different equations. The standard method
used in the sequential algorithm is a cell-centered finite volume method. Further schemes,
so far only available for the two-phase pressure equation, are cell-centered finite
volumes with multi-point flux approximation (Mpfa-O method) and mimetic finite differences.
An $h$-adaptive implementation of both sequential algorithms is provided for two dimensions.


\bibliographystyle{plainnat}
\bibliography{dumux-handbook}
\end{document}
